% ============================================================
%  CAPÍTULO 10 – Análisis
% ============================================================
\chapter{Análisis}
\label{ch:analisis}

\section{Requisitos Funcionales (RF)}

Los requisitos funcionales describen las capacidades que el sistema debe proporcionar a sus usuarios:

\begin{longtable}{@{}p{1.5cm}p{4cm}p{8.5cm}@{}}
\toprule
\textbf{ID} & \textbf{Módulo} & \textbf{Descripción} \\
\midrule
\endhead
RF-01 & Autenticación & El sistema permitirá el registro de nuevos usuarios mediante email, contraseña, nombre completo, fecha de nacimiento y teléfono. \\
RF-02 & Autenticación & El sistema permitirá el inicio de sesión mediante email y contraseña. \\
RF-03 & Autenticación & El sistema distinguirá entre usuarios con rol de \textit{entrenador} y rol de \textit{cliente}. \\
RF-04 & Resumen del Día & La aplicación mostrará al entrenador un resumen diario con el número de sesiones realizadas, pendientes y mensajes no leídos. \\
RF-05 & Resumen del Día & Se mostrarán los próximos entrenamentos del día con el cliente asignado. \\
RF-06 & Mis Clientes & El entrenador podrá consultar la lista de sus clientes con acceso a búsqueda por nombre. \\
RF-07 & Mis Clientes & El entrenador podrá acceder al perfil detallado de cada cliente (IMC, estado, sesiones realizadas, progreso de peso). \\
RF-08 & Mis Clientes & Se mostrarán las sesiones pasadas del cliente con valoración (sistema de estrellas). \\
RF-09 & Crear Rutina & El entrenador podrá crear una rutina asignada a un cliente, con título y ejercicios configurables. \\
RF-10 & Crear Rutina & El sistema permitirá buscar ejercicios en un catálogo local por nombre o grupo muscular, tras haber evaluado una API externa durante el desarrollo. \\
RF-11 & Crear Rutina & Cada ejercicio añadido a la rutina tendrá series, repeticiones y peso configurables. \\
RF-12 & Calendario & El entrenador podrá visualizar las sesiones programadas en un calendario mensual. \\
RF-13 & Calendario & Las sesiones se mostrarán con el nombre del cliente y el tipo de sesión. \\
RF-14 & Mensajes & El entrenador podrá consultar la lista de conversaciones con sus clientes. \\
RF-15 & Mensajes & El sistema permitirá enviar y recibir mensajes de texto en tiempo real. \\
RF-16 & Mensajes & Los mensajes no leídos se resaltarán con una notificación numérica. \\
\bottomrule
\caption{Requisitos funcionales de FitLabs.}
\label{tab:rf}
\end{longtable}

\section{Requisitos No Funcionales (RNF)}

\begin{longtable}{@{}p{1.5cm}p{3cm}p{9cm}@{}}
\toprule
\textbf{ID} & \textbf{Categoría} & \textbf{Descripción} \\
\midrule
\endhead
RNF-01 & Rendimiento & La aplicación debe responder a las interacciones del usuario en menos de 300 ms en condiciones de red normal. \\
RNF-02 & Seguridad & Todos los datos en tránsito deben transmitirse cifrados mediante TLS 1.2 o superior. \\
RNF-03 & Seguridad & Las políticas RLS de Supabase deben garantizar el aislamiento total de datos entre usuarios distintos. \\
RNF-04 & Disponibilidad & El backend en Supabase (plan Free) ofrece una disponibilidad del 99\,\% según SLA del proveedor. \\
RNF-05 & Usabilidad & La interfaz debe ser intuitiva, con una curva de aprendizaje asumible sin manual para un usuario con formación básica en smartphones. \\
RNF-06 & Escalabilidad & La arquitectura debe permitir incorporar nuevos módulos sin refactorización mayor del código existente. \\
RNF-07 & Compatibilidad & La aplicación debe ser compatible con Android 5.0 (API Level 21) o superior. \\
RNF-08 & Mantenibilidad & El código debe organizarse en módulos (pantallas) independientes, con navegación centralizada en \texttt{main.dart}. \\
\bottomrule
\caption{Requisitos no funcionales de FitLabs.}
\label{tab:rnf}
\end{longtable}

\section{Historias de usuario y criterios de aceptación}

Para convertir los requisitos en funcionalidades implementables, se redactaron historias de usuario en formato ágil. Esta capa intermedia permitió pasar de una definición general a criterios observables y verificables en pruebas.

\begin{longtable}{@{}p{1.3cm}p{5.5cm}p{6.7cm}@{}}
\toprule
\textbf{ID} & \textbf{Historia de usuario} & \textbf{Criterios de aceptación} \\
\midrule
\endhead
HU-01 & Como entrenador, quiero registrarme con mis datos para acceder a la app. & El alta crea usuario en Auth y perfil en BD; validación de campos obligatorios; mensaje de error claro si el email ya existe. \\
HU-02 & Como entrenador, quiero iniciar sesión para acceder a mi panel de trabajo. & Redirección a Resumen del Día con credenciales válidas; error controlado si la autenticación falla. \\
HU-03 & Como entrenador, quiero ver mis clientes para organizar el trabajo diario. & Solo se muestran clientes asociados al usuario autenticado; listado ordenado y cargado sin errores. \\
HU-04 & Como entrenador, quiero buscar clientes por nombre para ahorrar tiempo. & Filtro en tiempo real sin recarga completa de pantalla; admite coincidencias parciales. \\
HU-05 & Como entrenador, quiero crear rutinas personalizadas para cada cliente. & Guardado correcto de rutina y ejercicios vinculados; validación de título obligatorio. \\
HU-06 & Como entrenador, quiero consultar el calendario para planificar sesiones. & Días con sesiones visibles; detalle de sesiones al seleccionar fecha. \\
HU-07 & Como entrenador, quiero enviar mensajes para resolver dudas de clientes. & Mensaje persistido en BD y visualizado por ambos usuarios; actualización en tiempo real. \\
HU-08 & Como entrenador, quiero identificar mensajes no leídos para priorizar respuestas. & Contador de no leídos visible y actualizado tras lectura. \\
\bottomrule
\caption{Historias de usuario priorizadas para la versión entregada.}
\end{longtable}

\section{Priorización funcional (MoSCoW)}

Se aplicó una priorización MoSCoW para gestionar alcance y proteger la entrega final frente a cambios de última hora.

\begin{longtable}{@{}p{2cm}p{11.5cm}@{}}
\toprule
\textbf{Nivel} & \textbf{Funcionalidades asociadas} \\
\midrule
\endhead
Must Have & Registro/login, roles, resumen diario, gestión de clientes, detalle de cliente, creación de rutinas, calendario, mensajería en tiempo real. \\
Should Have & Indicadores visuales avanzados, limpieza automática de estados de formulario, mejoras de UX en flujos de alta frecuencia. \\
Could Have & Integración de API externa de ejercicios en producción, perfil avanzado con preferencias, métricas extendidas. \\
Won't Have (v1) & Pagos integrados, videollamadas, modo offline completo, integración con wearables. \\
\bottomrule
\caption{Priorización de alcance mediante técnica MoSCoW.}
\end{longtable}

\section{Casos de uso principales}

A nivel de modelado funcional, los casos de uso centrales del sistema son los siguientes:

\begin{enumerate}
  \item \textbf{CU-01: Gestionar autenticación}. Actor principal: Entrenador. Resultado: sesión iniciada y perfil cargado.
  \item \textbf{CU-02: Consultar cartera de clientes}. Actor principal: Entrenador. Resultado: listado filtrable y acceso a detalle.
  \item \textbf{CU-03: Diseñar rutina personalizada}. Actor principal: Entrenador. Resultado: rutina persistida con ejercicios parametrizados.
  \item \textbf{CU-04: Planificar sesiones en calendario}. Actor principal: Entrenador. Resultado: agenda mensual consultable y accionable.
  \item \textbf{CU-05: Mantener comunicación con clientes}. Actor principal: Entrenador/Cliente. Resultado: intercambio de mensajes en tiempo real.
\end{enumerate}

\begin{figure}[H]
  \centering
  \resizebox{0.8\textwidth}{!}{\includegraphics{casos_uso.png}}
  \caption{Diagrama de casos de uso de FitLabs con actores (Entrenador y Cliente) y funcionalidades principales.}
  \label{fig:casos_uso}
\end{figure}

Estos casos de uso se tomaron como referencia para estructurar tanto el diseño de pantallas como la planificación de pruebas del capítulo de validación.

\section{Arquitectura del sistema}

FitLabs sigue una arquitectura cliente-servidor \textbf{en tres capas}:

\begin{enumerate}
  \item \textbf{Capa de presentación} (cliente): Aplicación Flutter en el dispositivo del usuario. Gestiona la UI y la lógica de navegación.
  \item \textbf{Capa de negocio / API}: Supabase actúa como capa intermedia proporcionando la API PostgREST (generada automáticamente sobre el esquema PostgreSQL) y la API de autenticación.
  \item \textbf{Capa de datos}: Base de datos PostgreSQL alojada en la infraestructura cloud de Supabase, con políticas RLS aplicadas.
\end{enumerate}

Adicionalmente, durante el desarrollo se evaluó una \textbf{API REST externa} para el catálogo de ejercicios. La versión final del proyecto utiliza un catálogo local de ejercicios para reducir dependencias externas y asegurar el cierre funcional de la aplicación.

\begin{figure}[H]
  \centering
  % Sustituir por un diagrama real exportado como imagen
  \includegraphics[width=0.85\textwidth]{diagrama_arquitectura.png}
  \caption{Diagrama de arquitectura del sistema FitLabs.}
  \label{fig:arquitectura}
\end{figure}

% Exportar el diagrama de arquitectura desde draw.io o Figma y guardarlo como imagenes/diagrama_arquitectura.png.

\section{Esquema de la base de datos}

La base de datos de FitLabs está implementada en PostgreSQL a través de Supabase. A continuación se describe cada tabla del esquema:

\subsection{Tabla \texttt{perfiles}}

Almacena los perfiles de todos los usuarios (entrenadores y clientes).

\begin{table}[H]
\centering
\begin{tabular}{@{}llll@{}}
\toprule
\textbf{Campo} & \textbf{Tipo} & \textbf{PK/FK} & \textbf{Descripción} \\
\midrule
id & UUID & PK & Referencia al usuario de Supabase Auth \\
role & VARCHAR(50) & -- & Rol del usuario: \textit{trainer} o \textit{client} \\
username & VARCHAR(255) & UNIQUE & Nombre de usuario visible en la app \\
avatar\_url & TEXT & -- & URL del avatar almacenado en Supabase Storage \\
bio & TEXT & -- & Descripción breve del perfil \\
created\_at & TIMESTAMPTZ & -- & Fecha de creación del perfil \\
\bottomrule
\end{tabular}
\caption{Estructura de la tabla \texttt{perfiles}.}
\end{table}

\subsection{Tabla \texttt{clientes\_entrenador}}

Relación muchos-a-muchos entre entrenadores y clientes.

\begin{table}[H]
\centering
\begin{tabular}{@{}llll@{}}
\toprule
\textbf{Campo} & \textbf{Tipo} & \textbf{PK/FK} & \textbf{Descripción} \\
\midrule
id & UUID & PK & Identificador único \\
trainer\_id & UUID & FK $\to$ perfiles & Perfil del entrenador \\
client\_id & UUID & FK $\to$ perfiles & Perfil del cliente \\
status & VARCHAR(50) & -- & Estado: \textit{active}, \textit{pending} \\
created\_at & TIMESTAMPTZ & -- & Fecha de creación de la relación \\
\bottomrule
\end{tabular}
\caption{Estructura de la tabla \texttt{clientes\_entrenador}.}
\end{table}

\subsection{Tabla \texttt{rutinas}}

Almacena las rutinas creadas por los entrenadores.

\begin{table}[H]
\centering
\begin{tabular}{@{}llll@{}}
\toprule
\textbf{Campo} & \textbf{Tipo} & \textbf{PK/FK} & \textbf{Descripción} \\
\midrule
id & UUID & PK & Identificador único de la rutina \\
creator\_id & UUID & FK $\to$ perfiles & Entrenador creador \\
assigned\_client\_id & UUID & FK $\to$ perfiles & Cliente asignado (puede ser nulo) \\
title & VARCHAR(255) & -- & Nombre de la rutina \\
description & TEXT & -- & Descripción de la rutina \\
is\_public & BOOLEAN & -- & Si la rutina es visible para otros \\
created\_at & TIMESTAMPTZ & -- & Fecha de creación \\
\bottomrule
\end{tabular}
\caption{Estructura de la tabla \texttt{rutinas}.}
\end{table}

\subsection{Tabla \texttt{ejercicios\_rutina}}

Almacena los ejercicios incluidos en cada rutina con sus parámetros de carga.

\begin{table}[H]
\centering
\begin{tabular}{@{}llll@{}}
\toprule
\textbf{Campo} & \textbf{Tipo} & \textbf{Descripción} \\
\midrule
id & UUID & Identificador único \\
id\_rutina & UUID & FK $\to$ rutinas \\
id\_ejercicio\_externo & TEXT & ID del ejercicio en la API externa \\
serie & INT & Número de series \\
repeticiones & INT & Repeticiones por serie \\
peso & DECIMAL(10,2) & Peso en kilogramos \\
duracion & INT & Duración en segundos (ejercicios de tiempo) \\
descanso & INT & Tiempo de descanso en segundos \\
orden & INT & Posición del ejercicio en la rutina \\
\bottomrule
\end{tabular}
\caption{Estructura de la tabla \texttt{ejercicios\_rutina}.}
\end{table}

\subsection{Tablas de comunicación: \texttt{chats} y \texttt{mensajes}}

La tabla \texttt{chats} representa un canal de comunicación entre dos usuarios. La tabla \texttt{mensajes} almacena cada mensaje enviado en un chat.

Las columnas clave de \texttt{chats} son: \texttt{id\_usuario1} y \texttt{id\_usuario2} (FK $\to$ perfiles).\\
Las columnas clave de \texttt{mensajes} son: \texttt{id\_chat} (FK $\to$ chats), \texttt{id\_remitente} (FK $\to$ perfiles), \texttt{contenido} (TEXT) y \texttt{leido} (BOOLEAN).

\subsection{Tablas de comunidad: \texttt{seguidores}, \texttt{likes\_rutina} y \texttt{comentarios\_rutina}}

Estas tablas replican funcionalidades de red social básicas sobre rutinas compartidas públicamente: seguidores entre usuarios, reacciones y comentarios en rutinas.

\subsection{Relaciones y cardinalidades}

Para que el diagrama ER refleje correctamente la lógica del sistema, las relaciones clave deben representarse con su cardinalidad y participación:

\begin{itemize}
  \item \texttt{perfiles} a \texttt{rutinas}: una relación 1:N, ya que un entrenador puede crear muchas rutinas y cada rutina pertenece a un único creador.
  \item \texttt{perfiles} a \texttt{clientes\_entrenador}: una relación N:M materializada mediante una tabla intermedia, con participación obligatoria de ambas partes en la relación activa.
  \item \texttt{rutinas} a \texttt{ejercicios\_rutina}: una relación 1:N, dado que una rutina puede contener múltiples ejercicios.
  \item \texttt{chats} a \texttt{mensajes}: una relación 1:N, porque cada chat agrupa múltiples mensajes.
  \item \texttt{perfiles} a \texttt{chats}: dos relaciones 1:N desde cada usuario participante, con participación total del chat respecto a ambos usuarios.
\end{itemize}

Esta lectura ayuda a rehacer el diagrama de forma más precisa en draw.io, incorporando cardinalidad y participación explícitas para cada relación.

\clearpage
\begin{landscape}
\section{Diagrama Entidad-Relación}
\begin{figure}[H]
  \centering
  \includegraphics[width=0.98\linewidth,height=0.78\textheight,keepaspectratio]{diagrama_er.png}
  \caption{Diagrama Entidad-Relación del esquema de base de datos de FitLabs.}
  \label{fig:diagrama_er}
\end{figure}
\end{landscape}

\section{Diagrama de flujo -- Proceso de autenticación}

\begin{figure}[H]
  \centering
  \includegraphics[width=0.7\textwidth]{flujo_autenticacion.png}
  \caption{Diagrama de flujo del proceso de registro e inicio de sesión en FitLabs.}
  \label{fig:flujo_auth}
\end{figure}

% Crear flujo_autenticacion.png usando draw.io o similar con:
%  1. Usuario introduce email+contraseña
%  2. Llamada a Supabase Auth
%  3. Si éxito -> Crear perfil en tabla perfiles -> Redirigir a /resumen
%  4. Si error -> Mostrar mensaje de error

\subsection{Integración histórica con ExerciseDB}

Durante el desarrollo se llegó a probar una integración HTTP con ExerciseDB. El historial del repositorio conserva ese enfoque, que incluía paginación y timeout:

\begin{lstlisting}[caption={Fragmento histórico de consulta a ExerciseDB}]
final uri = Uri.https('exercisedb.dev', '/api/v1/exercises/search', {
  'q': query,
  'offset': offset.toString(),
  'limit': '10',
});
final response = await http.get(uri).timeout(const Duration(seconds: 10));
\end{lstlisting}

Ese intento se descartó al cambiar las condiciones de uso del proveedor, y la solución final quedó apoyada en un catálogo local de ejercicios.
